Estructura lipídica que establece el límite entre la célula y el espacio extracelular. Está formada por una capa doble de lípidos anfipáticos, que son llamados así por tener una zona polar y una no polar, dirigiendo la polar hacia dentro de la célula y la no polar hacia fuera.

\section{Lípidos de la membrana}

\subsection{Fosfolípidos}

Glicerol
2 colas de ácidos grasos
Cabeza de grupo fosfato

\subsection{Colesterol}

Cuatro anillos de carbono fusionados

\section{Incrustaciones}

\subsection{Carbohidratos}

Solo por fuera
Se unen formando glucoproteínas y glucolípidos

\subsection{Proteínas}

Incrustadas parcial o totalmente en ambos lados o a través de la membrana.
Pueden ser integrales o periféricas